\documentclass{beamer}
\mode<presentation>
\usepackage{amsmath}
\usepackage{amssymb}
%\usepackage{advdate}
\usepackage{adjustbox}
\usepackage{subcaption}
\usepackage{enumitem}
\usepackage{multicol}
\usepackage{mathtools}
\usepackage{listings}
\usepackage{physics}
\usepackage{url}
\def\UrlBreaks{\do\/\do-}
\usetheme{Boadilla}
\usecolortheme{lily}
\setbeamertemplate{footline}
{
  \leavevmode%
  \hbox{%
  \begin{beamercolorbox}[wd=\paperwidth,ht=2.25ex,dp=1ex,right]{author in head/foot}%
    \insertframenumber{} / \inserttotalframenumber\hspace*{2ex} 
  \end{beamercolorbox}}%
  \vskip0pt%
}

\providecommand{\nCr}[2]{\,^{#1}C_{#2}} % nCr
\providecommand{\nPr}[2]{\,^{#1}P_{#2}} % nPr
\providecommand{\mbf}{\mathbf}
\providecommand{\pr}[1]{\ensuremath{\Pr\left(#1\right)}}
\providecommand{\qfunc}[1]{\ensuremath{Q\left(#1\right)}}
\providecommand{\sbrak}[1]{\ensuremath{{}\left[#1\right]}}
\providecommand{\lsbrak}[1]{\ensuremath{{}\left[#1\right.}}
\providecommand{\rsbrak}[1]{\ensuremath{{}\left.#1\right]}}
\providecommand{\brak}[1]{\ensuremath{\left(#1\right)}}
\providecommand{\lbrak}[1]{\ensuremath{\left(#1\right.}}
\providecommand{\rbrak}[1]{\ensuremath{\left.#1\right)}}
\providecommand{\cbrak}[1]{\ensuremath{\left\{#1\right\}}}
\providecommand{\lcbrak}[1]{\ensuremath{\left\{#1\right.}}
\providecommand{\rcbrak}[1]{\ensuremath{\left.#1\right\}}}
\theoremstyle{remark}
\newtheorem{rem}{Remark}
\newcommand{\sgn}{\mathop{\mathrm{sgn}}}
\providecommand{\abs}[1]{\left\vert#1\right\vert}
\providecommand{\res}[1]{\Res\displaylimits_{#1}} 
\providecommand{\norm}[1]{\lVert#1\rVert}
\providecommand{\mtx}[1]{\mathbf{#1}}
\providecommand{\mean}[1]{E\left[ #1 \right]}
\providecommand{\fourier}{\overset{\mathcal{F}}{ \rightleftharpoons}}
%\providecommand{\hilbert}{\overset{\mathcal{H}}{ \rightleftharpoons}}
\providecommand{\system}{\overset{\mathcal{H}}{ \longleftrightarrow}}
	%\newcommand{\solution}[2]{\textbf{Solution:}{#1}}
%\newcommand{\solution}{\noindent \textbf{Solution: }}
\providecommand{\dec}[2]{\ensuremath{\overset{#1}{\underset{#2}{\gtrless}}}}
\newcommand{\myvec}[1]{\ensuremath{\begin{pmatrix}#1\end{pmatrix}}}
\let\vec\mathbf

\lstset{
%language=C,
frame=single, 
breaklines=true,
columns=fullflexible
}

\numberwithin{equation}{section}

\lstset{
  language=Python,
  basicstyle=\ttfamily\small,
  keywordstyle=\color{blue},
  stringstyle=\color{orange},
  numbers=left,
  numberstyle=\tiny\color{gray},
  breaklines=true,
  showstringspaces=false
}

\title{Problem 2.10.52}
\author{ee25btech11023-Venkata Sai}

\date{\today} 
\begin{document}

\begin{frame}
\titlepage
\end{frame}

\section*{Outline}
\begin{frame}
\tableofcontents
\end{frame}
\section{Problem}
\begin{frame}
\frametitle{Problem Statement}
%
Let $\vec{a} = \hat{\vec{i}} + 2\hat{\vec{j}} + \hat{\vec{k}}$, $\vec{b} = \hat{\vec{i}} - \hat{\vec{j}} + \hat{\vec{k}}$ and $\vec{c} = \hat{\vec{i}} + \hat{\vec{j}} - \hat{\vec{k}}$. A vector in the plane of $\vec{a}$ and $\vec{b}$ whose projection on $\vec{c}$ is $\frac{1}{\sqrt{3}}$, is 
\end{frame}
%\subsection{Literature}
\section{Solution}
\subsection{Formula}
\begin{frame}
\frametitle{Formula}
\setcounter{section}{1}
%\framesubtitle{Literature}
Let $\vec{v}$ be the vector in the plane of $\vec{a}$ and $\vec{b}$
\begin{align}
 [\vec{v} \quad \vec{a} \quad \vec{b}]=0 \implies \vec{v}=\alpha\vec {a}+\beta\vec {b} 
 \end{align}
 where
 \begin{align}
 \vec{a}=\myvec{1\\2\\1},\vec{b}=\myvec{1\\-1\\1}
 \end{align}
 Then
 \begin{align}
  \vec{v}=\alpha\myvec{1\\2\\1}+\beta\myvec{1\\-1\\1}
=\myvec{\alpha+\beta\\2\alpha-\beta\\\alpha+\beta}.
 \end{align}
\end{frame}
\subsection{Method of finding the projection}
\begin{frame}
\frametitle{Method of finding the projection}
Mathematically, the projection of $\vec{v}$ on $\vec{c}$ is defined as 
\begin{align}
\vec{d}=k\vec{c},\text{such that} \brak{\vec{v}-\vec{d}}^{\top}\vec{d}=0 
\end{align}
yielding
\begin{align}
\brak{\vec{v}-k\vec{c}}^{\top}\vec{d}=0 
\end{align}
\begin{align}
\text{or},
k=\frac{\vec{v}^{\top}\vec{c}}{\|\vec{c\|^{2}}}, \implies
\vec{d}=\frac{\vec v^{\top}\vec{c}}{\|\vec{c}\|^{2}}\;\vec{c}.  
\end{align}
\end{frame}
\subsection{Finding projection}
\begin{frame}
\frametitle{Finding projection}
\begin{align}
\norm{\vec{c}}^2=\vec{c}^\top\vec{c}&=\myvec{1&1&-1}\myvec{1\\1\\-1} \\
&=1+1+1=3  \implies \norm{\vec{c}}=\sqrt{3}
\end{align}
\begin{align}
\vec{v}^{\top}\vec{c}&=\myvec{\alpha+\beta & 2\alpha-\beta & \alpha+\beta}\myvec{1\\1\\-1} \\
&=\alpha+\beta + 2\alpha-\beta - \alpha-\beta=2\alpha-\beta
\end{align}
\begin{align}
\vec{d}=\frac{2\alpha-\beta}{3}\vec{c}
\end{align}
\end{frame}

\subsection{Conclusion}
\begin{frame}
\frametitle{Conclusion}
Given $\norm{\vec{d}}=\frac{1}{\sqrt{3}}$ , Then 
\begin{align}
\frac{|2\alpha-\beta|}{3}\norm{\vec{c}}=\frac{|2\alpha-\beta|}{3}\sqrt{3}=\frac{1}{\sqrt{3}}  \implies |2\alpha - \beta| = 1 
\end{align}
From equation(3) we know that that $x$ and $z$ coordinates in our required vector must be equal. \\
If $2\alpha - \beta $ = -1,from the most reliable option i.e option(1) which has the same $x$ and $z$ coordinates, we get 
\begin{align}
    \alpha+\beta=4\ \text{and}\ 2\alpha-\beta=-1 
\end{align}
\begin{align}
    3\alpha=3 \implies \alpha =1\ \text{and}\ \beta=3 \implies \vec{v}=\myvec{4 \\ -1\\ 4}
    \end{align}
If $2\alpha -\beta $ =1,from the options $x$ and $z$ coordinates are
not equal,No option satisfy.
\end{frame}

\subsection{Plot}
\begin{frame}[fragile]
\frametitle{Plot}

\begin{figure}[h!]
   \centering
   \includegraphics[width=0.7\columnwidth]{figs/fig1.png}
	\caption{}
   \label{stemplot}
\end{figure}
\end{frame}

\section{C Code}
\begin{frame}[fragile]
\frametitle{C Code}
\begin{lstlisting}[language=C]
void get_vectors(double* vector_data) {
    vector_data[0] = 1.0; vector_data[1] = 2.0; vector_data[2] = 1.0;
    vector_data[3] = 1.0; vector_data[4] = -1.0; vector_data[5] = 1.0;
    vector_data[6] = 1.0; vector_data[7] = 1.0; vector_data[8] = -1.0;
    vector_data[9] = 4.0; vector_data[10] = -1.0; vector_data[11] = 4.0;
}
    \end{lstlisting}
\end{frame}
\section{Python Code}
\begin{frame}[fragile]
\frametitle{Calling C Function}
\begin{lstlisting}[language=Python]
import ctypes
import numpy as np
def load_vectors_from_c():
    lib = ctypes.CDLL('./plane.so')
    # Define the C function signature
    get_vectors_c = lib.get_vectors
    get_vectors_c.argtypes = [ctypes.POINTER(ctypes.c_double)]
    get_vectors_c.restype = None

    vector_data = np.zeros(12, dtype=np.float64)
    get_vectors_c(vector_data.ctypes.data_as(ctypes.POINTER(ctypes.c_double)))
    # Reshape the data and return the individual vectors
    vectors_from_c = vector_data.reshape(4, 3)
    a, b, c, v = vectors_from_c[0], vectors_from_c[1], vectors_from_c[2], vectors_from_c[3]
    return a, b, c, v
\end{lstlisting}
\end{frame}

\begin{frame}[fragile]
\frametitle{Python Code for Plotting}
\begin{lstlisting}[language=Python]
#Code by GVV Sharma
#September 12, 2023
#Revised July 21, 2024
#released under GNU GPL

import sys
import matplotlib.pyplot as plt
import numpy as np
sys.path.insert(0, '/workspaces/urban-potato/matgeo/codes/CoordGeo/') 
from call import load_vectors_from_c
hat_symbol = '\u0302'
from line.funcs import *
from triangle.funcs import *
from conics.funcs import circ_gen

a, b, c, v = load_vectors_from_c()
 
\end{lstlisting}
\end{frame}

\begin{frame}[fragile]
\frametitle{Python Code for Plotting}
\begin{lstlisting}[language=Python]
vx, vy, vz = v[0], v[1], v[2]
j_part_str = f"+ {vy}j{hat_symbol}" if vy >= 0 else f"- {abs(vy)}j{hat_symbol}"
k_part_str = f"+ {vz}k{hat_symbol}" if vz >= 0 else f"- {abs(vz)}k{hat_symbol}"
print(f"The vector in the plane of a and b is: {vx}i{hat_symbol} {j_part_str} {k_part_str}")

# Create a 3D plot
fig = plt.figure(figsize=(10, 8))
ax = fig.add_subplot(111, projection='3d')
origin = [0, 0, 0]

# Plot vectors
vectors = {'a': (a, 'r'), 'b': (b, 'g'), 'c': (c, 'b'), 'v': (v, 'orange')}
for name, (vec, color) in vectors.items():
    x, y, z = int(vec[0]), int(vec[1]), int(vec[2])
    j_part = f"+ {y}j{hat_symbol}" if y >= 0 else f"- {abs(y)}j{hat_symbol}"
\end{lstlisting}
\end{frame}

\begin{frame}[fragile]
\frametitle{Python Code for Plotting}
\begin{lstlisting}[language=Python]
    k_part = f"+ {z}k{hat_symbol}" if z >= 0 else f"- {abs(z)}k{hat_symbol}"
    label_text = f'${name}$ = {x}i{hat_symbol} {j_part} {k_part}'
  ax.quiver(origin[0], origin[1], origin[2], vec[0], vec[1], vec[2],
              color=color, label=label_text,
              arrow_length_ratio=0.1, linewidth=2)
    ax.text(vec[0]*1.1, vec[1]*1.1, vec[2]*1.1, f'${name}$', color=color, fontsize=15)
    # Plot the plane
x_plane = np.linspace(-5, 5, 10); y_plane = np.linspace(-5, 5, 10)
X, Y = np.meshgrid(x_plane, y_plane)
Z = X
ax.plot_surface(X, Y, Z, alpha=0.2, color='gray')

\end{lstlisting}
\end{frame}

\begin{frame}[fragile]
\frametitle{Python Code for Plotting}
\begin{lstlisting}[language=Python]
ax.view_init(elev=25, azim=45)
ax.set_xlabel('X'); ax.set_ylabel('Y'); ax.set_zlabel('Z')
ax.set_title('Finding a vector that is in the plane of vectors a and b')
ax.legend()
ax.grid(True)
ax.set_xlim([-5, 5]); ax.set_ylim([-5, 5]); ax.set_zlim([-5, 5])
\end{lstlisting}
\end{frame}

\end{document}
